\subsection{\texorpdfstring{অক্টাল থেকে $n$ ভিত্তি}{অক্টাল থেকে n ভিত্তি}}
প্রথমে স্থানমান ব্যবহার করে দশমিকে আনো; তারপর লক্ষ্য ভিত্তিতে ভাগ ও গুণ করো। দুই ধাপের পূর্ণাংশ ও ভগ্নাংশ আলাদাভাবে মিলিয়ে নাও।
\par\noindent\begin{minipage}{\linewidth}
\convExample{১}{পূর্ণসংখ্যা}
$ (73)_{8} $-কে ৬ ভিত্তিতে রূপান্তর করো।\par
\textbf{প্রথম ধাপ: দশমিক মান}\par
\par\noindent\begin{minipage}{\linewidth}\centering
\phantomsection\label{VC042}
\begin{chThreeDiagram}\begin{tikzpicture}[x=1mm,y=1mm]
\node[vc digit] (d0) at (0,0) {$ 7 $};
\node[vcnum] (w0) at (0,-10) {$ 8^{1} $};
\node[vcnum] (v0) at (0,-25) {$ 56 $};
\draw[vc arrow] (d0.south)--(w0.north);
\draw[vc arrow] (w0.south)--(v0.north);
\node[vc digit] (d1) at (19,0) {$ 3 $};
\node[vcnum] (w1) at (19,-10) {$ 8^{0} $};
\node[vcnum] (v1) at (19,-25) {$ 3 $};
\draw[vc arrow] (d1.south)--(w1.north);
\draw[vc arrow] (w1.south)--(v1.north);
\end{tikzpicture}\end{chThreeDiagram}
\end{minipage}\par

\[\begin{aligned}V&=7\times 8^{1}+3\times 8^{0}\\&=56+3=59.\end{aligned}\]
\textbf{দ্বিতীয় ধাপ: লক্ষ্য ভিত্তি}\par
\par\noindent\begin{minipage}{\linewidth}\centering
\phantomsection\label{VC043}
\begin{chThreeDiagram}\begin{tikzpicture}[x=1mm,y=1mm]
\path[use as bounding box] (-3,9) rectangle (72,-35);
\node[vcnote] at (15,7) {ভাজ্য/ভাগফল};
\node[vcnote] at (44,7) {ভাগশেষ};
\node[vcnum] at (0.0,0) {$ 6 $};
\node[vcnum] at (10.0,0) {$ 59 $};
\draw[vc line] (5.0,3)--(5.0,-3.5)--(50,-3.5);
\node[vcnum] (r0) at (44,-7) {$ 5 $};
\node[vcnum] at (1.5,-7) {$ 6 $};
\node[vcnum] at (11.5,-7) {$ 9 $};
\draw[vc line] (6.5,-4)--(6.5,-10.5)--(50,-10.5);
\node[vcnum] (r1) at (44,-14) {$ 3 $};
\node[vcnum] at (3.0,-14) {$ 6 $};
\node[vcnum] at (13.0,-14) {$ 1 $};
\draw[vc line] (8.0,-11)--(8.0,-17.5)--(50,-17.5);
\node[vcnum] (r2) at (44,-21) {$ 1 $};
\node[vcnum] at (14.5,-21) {$0$};
\node[vcnote] at (61,-3) {LSB};
\node[vcnote] at (61,-26) {MSB};
\draw[vc arrow] (61,-21)--(61,-7);
\node[vcnote] at (31,-33) {নিচ থেকে ওপরে পড়ো};
\end{tikzpicture}\end{chThreeDiagram}
\end{minipage}\par
\[\therefore\ (73)_{8}=(135)_{6}.\]
\end{minipage}\par
\par\noindent\begin{minipage}{\linewidth}
\convExample{২}{শুধু ভগ্নাংশ}
$ (0.4)_{8} $-কে ৬ ভিত্তিতে রূপান্তর করো।\par
\textbf{প্রথম ধাপ: দশমিক মান}\par
\par\noindent\begin{minipage}{\linewidth}\centering
\phantomsection\label{VC044}
\begin{chThreeDiagram}\begin{tikzpicture}[x=1mm,y=1mm]
\node[vc digit] (d0) at (0,0) {$ 0 $};
\node[vcnum] (w0) at (0,-10) {$ 8^{0} $};
\node[vcnum] (v0) at (0,-25) {$ 0 $};
\draw[vc arrow] (d0.south)--(w0.north);
\draw[vc arrow] (w0.south)--(v0.north);
\node[vc digit] (d1) at (27,0) {$ 4 $};
\node[vcnum] (w1) at (27,-10) {$ 8^{-1} $};
\node[vcnum] (v1) at (27,-25) {$ \frac{1}{2} $};
\draw[vc arrow] (d1.south)--(w1.north);
\draw[vc arrow] (w1.south)--(v1.north);
\node[vc digit] at (13.5,0) {$.$};
\end{tikzpicture}\end{chThreeDiagram}
\end{minipage}\par

\[\begin{aligned}V&=0\times 8^{0}+4\times 8^{-1}\\&=\frac{1}{2}=0.5.\end{aligned}\]
\textbf{দ্বিতীয় ধাপ: লক্ষ্য ভিত্তি}\par
\par\noindent\begin{minipage}{\linewidth}\centering
\phantomsection\label{VC045}
\begin{chThreeDiagram}\begin{tikzpicture}[x=1mm,y=1mm]
\path[use as bounding box] (-9,9) rectangle (76,-27);
\node[vcnote] at (14,7) {প্রাপ্ত অঙ্ক};
\node[vcnote] at (45,7) {ভগ্নাংশ};
\node[vcnum] at (45,0) {$ 0.5 $};
\draw[vc line] (26,3)--(26,-15);
\node[vcnum] (m0) at (45,-5) {$\times 6$};
\draw[vc line] (7,-8)--(65,-8);
\node[vcnum] at (14,-12) {$ 3 $};
\node[vcnum] (f0) at (45,-12) {$ 0 $};
\node[vcnote] at (0,-10) {MSB=LSB};
\draw[vc arrow] (0,-14)--(0,-20);
\node[vcnote] at (34,-24) {ওপর থেকে নিচে পড়ো};
\end{tikzpicture}\end{chThreeDiagram}
\end{minipage}\par
\[\therefore\ (0.4)_{8}=(0.3)_{6}.\]
\end{minipage}\par
\par\noindent\begin{minipage}{\linewidth}
\convExample{৩}{মিশ্র সংখ্যা}
$ (125.6)_{8} $-কে ৬ ভিত্তিতে রূপান্তর করো।\par
\textbf{প্রথম ধাপ: দশমিক মান}\par
\par\noindent\begin{minipage}{\linewidth}\centering
\phantomsection\label{VC046}
\begin{chThreeDiagram}\begin{tikzpicture}[x=1mm,y=1mm]
\node[vc digit] (d0) at (0,0) {$ 1 $};
\node[vcnum] (w0) at (0,-10) {$ 8^{2} $};
\node[vcnum] (v0) at (0,-25) {$ 64 $};
\draw[vc arrow] (d0.south)--(w0.north);
\draw[vc arrow] (w0.south)--(v0.north);
\node[vc digit] (d1) at (19,0) {$ 2 $};
\node[vcnum] (w1) at (19,-10) {$ 8^{1} $};
\node[vcnum] (v1) at (19,-25) {$ 16 $};
\draw[vc arrow] (d1.south)--(w1.north);
\draw[vc arrow] (w1.south)--(v1.north);
\node[vc digit] (d2) at (38,0) {$ 5 $};
\node[vcnum] (w2) at (38,-10) {$ 8^{0} $};
\node[vcnum] (v2) at (38,-25) {$ 5 $};
\draw[vc arrow] (d2.south)--(w2.north);
\draw[vc arrow] (w2.south)--(v2.north);
\node[vc digit] (d3) at (65,0) {$ 6 $};
\node[vcnum] (w3) at (65,-10) {$ 8^{-1} $};
\node[vcnum] (v3) at (65,-25) {$ \frac{3}{4} $};
\draw[vc arrow] (d3.south)--(w3.north);
\draw[vc arrow] (w3.south)--(v3.north);
\node[vc digit] at (51.5,0) {$.$};
\end{tikzpicture}\end{chThreeDiagram}
\end{minipage}\par

\[\begin{aligned}V&=1\times 8^{2}+2\times 8^{1}+5\times 8^{0}+6\times 8^{-1}\\&=64+16+5+\frac{3}{4}=85.75.\end{aligned}\]
\textbf{দ্বিতীয় ধাপ: লক্ষ্য ভিত্তি}\par
\begin{minipage}[t]{.48\linewidth}\vspace{0pt}\centering\textbf{পূর্ণাংশ}\par\vspace{1mm}
\par\noindent\begin{minipage}{\linewidth}\centering
\phantomsection\label{VC047}
\begin{chThreeDiagram}\begin{tikzpicture}[x=1mm,y=1mm]
\path[use as bounding box] (-3,9) rectangle (72,-35);
\node[vcnote] at (15,7) {ভাজ্য/ভাগফল};
\node[vcnote] at (44,7) {ভাগশেষ};
\node[vcnum] at (0.0,0) {$ 6 $};
\node[vcnum] at (10.0,0) {$ 85 $};
\draw[vc line] (5.0,3)--(5.0,-3.5)--(50,-3.5);
\node[vcnum] (r0) at (44,-7) {$ 1 $};
\node[vcnum] at (1.5,-7) {$ 6 $};
\node[vcnum] at (11.5,-7) {$ 14 $};
\draw[vc line] (6.5,-4)--(6.5,-10.5)--(50,-10.5);
\node[vcnum] (r1) at (44,-14) {$ 2 $};
\node[vcnum] at (3.0,-14) {$ 6 $};
\node[vcnum] at (13.0,-14) {$ 2 $};
\draw[vc line] (8.0,-11)--(8.0,-17.5)--(50,-17.5);
\node[vcnum] (r2) at (44,-21) {$ 2 $};
\node[vcnum] at (14.5,-21) {$0$};
\node[vcnote] at (61,-3) {LSB};
\node[vcnote] at (61,-26) {MSB};
\draw[vc arrow] (61,-21)--(61,-7);
\node[vcnote] at (31,-33) {নিচ থেকে ওপরে পড়ো};
\end{tikzpicture}\end{chThreeDiagram}
\end{minipage}\par

$ (221)_{6} $
\end{minipage}\hfill
\begin{minipage}[t]{.48\linewidth}\vspace{0pt}\centering\textbf{ভগ্নাংশ}\par\vspace{1mm}
\par\noindent\begin{minipage}{\linewidth}\centering
\phantomsection\label{VC048}
\begin{chThreeDiagram}\begin{tikzpicture}[x=1mm,y=1mm]
\path[use as bounding box] (-9,9) rectangle (76,-39);
\node[vcnote] at (14,7) {প্রাপ্ত অঙ্ক};
\node[vcnote] at (45,7) {ভগ্নাংশ};
\node[vcnum] at (45,0) {$ 0.75 $};
\draw[vc line] (26,3)--(26,-27);
\node[vcnum] (m0) at (45,-5) {$\times 6$};
\draw[vc line] (7,-8)--(65,-8);
\node[vcnum] at (14,-12) {$ 4 $};
\node[vcnum] (f0) at (45,-12) {$ 0.5 $};
\node[vcnum] (m1) at (45,-17) {$\times 6$};
\draw[vc line] (7,-20)--(65,-20);
\node[vcnum] at (14,-24) {$ 3 $};
\node[vcnum] (f1) at (45,-24) {$ 0 $};
\draw[vc arrow] (f0.east)--(72,-12)--(72,-17)--(m1.east);
\node[vcnote] at (0,-8) {MSB};
\node[vcnote] at (0,-29) {LSB};
\draw[vc arrow] (0,-12)--(0,-24);
\node[vcnote] at (34,-36) {ওপর থেকে নিচে পড়ো};
\end{tikzpicture}\end{chThreeDiagram}
\end{minipage}\par

$ (0.43)_{6} $
\end{minipage}\par
\[\therefore\ (125.6)_{8}=(221.43)_{6}.\]
\end{minipage}\par
